\documentclass[../veristore_security_model.tex]{subfiles}

\begin{document}
\section{Verification and Testing Plan}

% COMMENT: Outline how you will verify that security mitigations are implemented correctly and
% continue to work as the system evolves. This includes unit tests for security-critical code,
% integration tests for security protocols, penetration testing, code review process, and
% regression testing. Specify who is responsible for creating and maintaining security tests.

\subsection{Security Test Suite}

% COMMENT: Describe the automated test suite that will verify security properties. Tests should
% cover: fingerprint collision resistance, detection of corrupted fragments, verification protocol
% correctness, proper handling of malformed input. Each mitigated threat should have corresponding
% tests.

\subsection{Code Review Process}

In review, we check four things:
\begin{itemize}
    \item The change does not weaken integrity checks or bypass verification paths.
    \item Error handling is safe and returns clear HTTP status codes.
    \item Input validation is still strict for API payloads.
    \item Tests were added or updated for the changed behavior.
\end{itemize}

\subsection{Regression Testing}

We keep a small set of repeatable security checks:
\begin{itemize}
    \item Byzantine fragment detection still rejects corrupted data.
    \item Retrieval still succeeds with only $m$ healthy servers.
    \item Edge-case payload tests still pass, including binary and large inputs.
\end{itemize}

\newpage
\end{document}
